% ============================================================
% Section 3.3 — System Identification Framework
% Target: ~0.5 page of 2-column IEEE conference format
% Status: 2nd DRAFT 2026-05-15 (label refs verified after §3.1↔§3.2 swap;
%         alg:ssa-jacobi still resolves — algorithm now lives in §3.2,
%         label is content-anchored)
% ============================================================

\subsection{System Identification Framework}
\label{sec:sysid-framework}

% --- Paragraph 1: Loss design per signal type ---
System identification treats Algorithm~\ref{alg:ssa-jacobi} as a
differentiable forward map
$\theta \mapsto \mathbf{y}_{\text{sim}}(\theta)$ and uses Adam to fit
measured signals $\mathbf{y}_{\text{real}}$. We calibrate parameters
sequentially in two stages with parameter-specific loss design.

\textbf{Stage 1 (elastic, one-step).} For the bulk stiffness $k$,
each Adam step uses a one-step loading experiment from the rest
state: a blunt indenter is pushed into the mesh and the loss is the
MSE between simulated and measured knife reaction force versus
displacement, in the spirit of Kolawole's volume-MSE inverse
FEM~\cite{kolawole2023} with displacement as the boundary input and
force as the target. We emphasize that the recovered $k$ is the
\emph{effective stiffness at our chosen $(K, dt)$}: shape-matching
PBD's effective compliance depends on iteration count and substep
size~\cite{macklin2016xpbd}, so $k$ should not be transferred to
solver configurations with different $(K, dt)$ without
re-calibration.

\textbf{Stage 2 (cutting, trajectory replay).} For the cutting
parameters---the $P(F)$ sharpness $\kappa$ (\ref{eq:pf}), the
participation-mask sharpness $s$, and the damage threshold $\tau$
(\ref{eq:soft-removal})---we hold $k$ fixed at the Stage-1 value and
replay a prescribed knife trajectory from the rest configuration.
The loss is the per-vertex consistency term
$\mathcal{L} \!=\! \tfrac{1}{|\mathcal{V}|}\sum_v\big(\alpha_v -
(1 - \phi_{\text{phantom}}(\mathbf{p}_v))\big)^2$,
where $\phi_{\text{phantom}}: \mathbb{R}^3 \to [0,1]$ is a smooth
cavity mask derived from the phantom's reconstructed signed-distance
field. Endpoint-only supervision (loss evaluated only on the final
cavity) avoids per-step accumulation but, unlike the Stage-1
protocol, gradients still backpropagate through the full multi-substep
trajectory. We accept this long-horizon backprop as an inherent cost
of cutting calibration; in our experiments the horizon is short
enough ($\leq\!10$ substeps per Adam step) that the
bias-variance tradeoff identified by Suh~\textit{et~al.}~\cite{suh2022}
remains acceptable.
Two descent channels contribute: a direct
$\partial\mathcal{L}/\partial\alpha_v$ path and an indirect
$\partial\mathcal{L}/\partial\mathbf{p}_v$ path that propagates
through the cluster solve via (\ref{eq:bridge}), exploiting the
position dependence to constrain cutting parameters beyond pure mask
matching. Because $\kappa$ and $s$ both control cavity-boundary
sharpness through different sigmoids, their joint identifiability is
not obvious from the loss alone; we report the empirical
identifiability landscape (Hessian conditioning and pairwise
parameter correlation) in \S\ref{sec:sysid-synth}.

% --- Paragraph 2: 3-vs-2 decision matrix ---
Friction $\mu$ and damping $\beta$, by contrast, are not identifiable
under either Stage-1 or Stage-2 protocols because their signal
requires nonzero particle velocity. They instead require a dedicated
stress-relaxation or free-vibration experiment with a time-series
force loss; when such hardware is unavailable, we adopt literature
values ($\mu$ from Misra~\textit{et~al.}~\cite{misra2010}, $\beta$
from Mackle~\textit{et~al.}~\cite{mackle2020}), consistent with the
convention in DiSECt~\cite{heiden2021disect} and the cardiac inverse
FEM work of Kolawole~\cite{kolawole2023}. Recovery accuracy and
convergence behaviour for $k$, $\kappa$, $s$, and $\tau$ under the
above protocols are evaluated in \S\ref{sec:sysid-synth}.
